\section{Configura\c{c}\~ao Experimental}
\label{sec:setup}

Os experimentos respondem a duas perguntas: se a sentença que a UDV aponta como evidência sustenta a
opinião, o que é medido por anotação humana cega numa amostra das audiências de teste
(Seção~\ref{sec:setup-udv}), e se o perfil e os trechos recuperados aumentam a capacidade do modelo
simulador de reconhecer a opinião que a matéria atribuiu à pessoa numa audiência que não entrou no
perfil (Seção~\ref{sec:setup-simulacao}). Na segunda avaliação, a UDV apenas liga a opinião a um
ator, pelo turno da evidência, e a sentença não aparece na pergunta. O resultado da simulação depende,
portanto, da resolução da pessoa na audiência e da identificação do ator entre audiências
(Seções~\ref{sec:metodo-udv} e~\ref{sec:metodo-atores}), sem depender da precisão do vínculo por
similaridade semântica que a primeira avaliação mede.

\subsection{Validação humana da evidência}
\label{sec:setup-udv}

\input{tables/validation-sample}

% Pendente: o artigo ainda não relata a comparação do codificador com as alternativas de recuperação nas audiências de validação; opção: uma frase no parágrafo dos controles da Seção 6.1 (mesmo teste de 142/40 consultas), a partir de retrieval_v1_report.json, com entrada bib para a correção de Holm.
A precisão da evidência é medida numa amostra cega sorteada somente nas audiências de teste, porque as
de validação estavam reservadas, antes do sorteio, à comparação do codificador com alternativas de
recuperação da evidência, e uma estimativa tirada delas tenderia a favorecer a alternativa escolhida.
A amostra é estratificada pelo tipo de vínculo (Tabela~\ref{tab:amostra}), e as UDVs sem evidência
de uma mesma pessoa ocupam uma só linha, porque a pergunta feita a elas depende só da pessoa.

Entre as UDVs semânticas, a amostra separa as que têm \emph{corroboração fraca}, marca que a UDV
registra, sem alterar o nível, quando o prefixo mais longo da citação encontrado na fala da pessoa tem
menos de seis palavras, o que não basta para o nível \texttt{quote\_found}, e a primeira ocorrência
desse prefixo está na sentença escolhida pelo codificador. Nessas linhas, como nas de citação literal,
as palavras entre aspas reaparecem na sentença, que tem assim um apoio textual além do cosseno. Com a
separação, o estrato \texttt{semantic\_match\_high}, avaliado por um dos critérios, reúne apenas
vínculos por similaridade.

Nas linhas com evidência, o anotador vê o nome e o cargo da pessoa, a opinião, a sentença e o texto ao
redor, e pode consultar o turno na transcrição, mas não vê o nível, o cosseno nem a marca de
corroboração fraca. Ele julga a sentença como \emph{correta}, quando sustenta o conteúdo principal da
opinião, inclusive a posição e os números ou nomes citados, \emph{parcial}, quando sustenta só parte da
opinião, a mesma posição sem o detalhe principal ou tem o apoio apenas no texto ao redor, ou
\emph{incorreta}, quando trata de outro assunto, tem a posição oposta, é condução da sessão ou é fala
de outra pessoa, o que ele confere pelo cabeçalho do turno. Como as palavras citadas reaparecem na
sentença nas linhas de citação literal e de corroboração fraca, o anotador pode reconhecer esses
estratos, e só a distinção entre \texttt{semantic\_match\_high} e \texttt{semantic\_match\_weak}
permanece cega. Nas linhas sem evidência, o anotador indica se a pessoa falou na sessão, e a UDV está
correta quando ela não falou.
% Decisão pendente: a repetição de 20 linhas depois de 24 horas não tem artefato; volta aqui só se for feita antes da submissão.

Os critérios, registrados antes de qualquer julgamento, usam o limite inferior do intervalo de
Wilson~\citep{wilson1927probable} de 95\% da precisão estrita (só ``correta'') na citação literal e da
precisão tolerante (``correta'' ou ``parcial'') no estrato \texttt{semantic\_match\_high}; os limiares
e o número de acertos que eles exigem com os tamanhos sorteados estão na Tabela~\ref{tab:amostra}.

\subsection{Avaliação da simulação por múltipla escolha}
\label{sec:setup-simulacao}

\input{figures/simulation}

\textbf{Perguntas.} Cada UDV de uma audiência de validação ou de teste cujo turno de evidência
pertence a um ator com perfil gera uma pergunta, acompanhada da data e do assunto da audiência
(Figura~\ref{fig:simulacao}). O texto da pergunta é ``Uma destas afirmações foi feita por
\textit{nome} nesta audiência, que não aparece no material desta conversa. Qual é a mais provável?''.
As quatro opções são a opinião da UDV, na redação do resumo estruturado, e uma opinião de cada um de
três outros participantes da mesma audiência, com participantes e opiniões sorteados com semente fixa.
As opiniões da própria pessoa e os textos iguais a elas são excluídos, um texto atribuído a mais de um
participante é incluído uma única vez, e a pergunta com menos de três outros participantes é
descartada. Como os distratores vêm da mesma audiência, eles tratam do mesmo tema, o que reduz, sem
eliminar, a possibilidade de escolher a opção pela proximidade com o assunto. Das 101 perguntas de
teste (Seção~\ref{sec:resultados-simulacao}), 84 vêm de UDVs \texttt{semantic\_match\_high}, 15 de
\texttt{quote\_found} e 2 de \texttt{semantic\_match\_weak}. Antes da execução, verifica-se que
nenhum perfil e nenhuma fala usada na recuperação contém audiência de validação ou de teste; a
verificação não cobre o cargo, que vem da matéria da audiência avaliada.

\textbf{Leitura da resposta.} A escolha é lida na distribuição de probabilidade do primeiro token
gerado, restrita aos tokens das quatro letras, e antes da execução confere-se que cada letra é um
único token. Como LLMs favorecem certas posições entre as
opções~\citep{zheng2024large,pezeshkpour2024large}, cada pergunta é apresentada nas quatro rotações
cíclicas de uma ordem sorteada. Em cada rotação, as probabilidades das letras são renormalizadas
entre si. A probabilidade de cada opção é a
média das rotações, a resposta é a opção de maior probabilidade, e as métricas são o acerto e a
probabilidade média da opção correta.

Como a leitura sempre escolhe uma letra e a probabilidade do primeiro token pode divergir da
resposta escrita~\citep{wang2024answer}, cada condição registra como diagnóstico a soma das
probabilidades das quatro letras no vocabulário inteiro. Essa soma mostra quanto da distribuição do
primeiro token cabe às letras.

\textbf{Seleção de $k$.} O número de trechos é, para cada modelo, o $k \in \{3, 5, 10\}$ de maior
acerto da condição 2 nas audiências de validação, com empates resolvidos pela maior probabilidade média da opção correta
e, depois, pelo menor valor da grade. No teste, a condição 2 é medida só com esse $k$.

\textbf{Comparação entre condições.} A diferença de acerto entre condições vizinhas estima o efeito
do elemento acrescentado, e o intervalo de 95\% de cada diferença é obtido por \textit{bootstrap}
pareado de percentis, com 1.000 reamostragens de audiências inteiras~\citep{efron1993introduction}.
As audiências são reamostradas inteiras porque as perguntas de uma audiência compartilham assunto,
falantes e distratores. Os quatro intervalos, dois por modelo, não são ajustados para comparações
múltiplas, e as demais medidas são descritivas.

\textbf{Modelos e recursos.} Avaliamos dois LLMs, o Gemma 4 31B e o Qwen3.8 27B, este com pesos
quantizados em 8 bits. Cada um escreve os 264 perfis a partir das falas de treino e responde às
perguntas como modelo simulador. Na geração dos perfis, os dois prompts
(Seção~\ref{sec:metodo-perfis}) compõem uma conversa de duas mensagens. Na geração e na simulação, as
mensagens são formatadas pelo \textit{chat template} do próprio LLM, sem formatação adicional e com o
modo de raciocínio desativado. A geração usa amostragem com temperatura 0,2, \textit{top-p} de 0,9 e
limite de 3.000 tokens de saída. O Gemma é executado em três GPUs A6000 com 48 GB de VRAM cada, em
bf16, e o Qwen, num Apple M4 Max com 128 GB de memória. Na recuperação dos trechos, o codificador é o
Serafim 335m, na mesma versão de pesos da UDV. O código, as configurações, os prompts e as saídas
estão em <repositório>.
